\chapter{Continuidad uniforme y control de las distancias}\label{chap:07}
\section{Una misma precisión para todos los puntos}
La continuidad permite elegir la precisión de entrada después de fijar un punto. La continuidad uniforme exige una elección que funcione simultáneamente en todo el dominio. Esta diferencia de cuantificadores explica por qué una función puede comportarse bien cerca de cada punto y, sin embargo, amplificar mucho las distancias en regiones cada vez más lejanas. Usaremos continuidad, sucesiones y la desigualdad triangular; no necesitaremos compacidad.

\begin{definicion}\label{ch07:def-uniforme}
Una aplicación $f:(X,d)\to(Y,\rho)$ es \emph{uniformemente continua} si
\[
 \forall\eps>0\ \exists\delta>0\ \forall x,y\in X,
 \qquad d(x,y)<\delta\Longrightarrow\rho(f(x),f(y))<\eps.
\]
Es \emph{Lipschitz} con constante $L\geq0$ si
\[
 \rho(f(x),f(y))\leq Ld(x,y)\qquad(x,y\in X).
\]
\end{definicion}
\index{continuidad!uniforme}\index{Lipschitz, funcion de@Lipschitz, función de}
La desigualdad de Lipschitz es no estricta: escribir $<$ para todo par de puntos produciría $0<0$ al tomar $x=y$. Una constante de Lipschitz no tiene por qué ser óptima; cualquier constante mayor también sirve.

\begin{proposicion}\label{ch07:lip-uniforme}
Toda aplicación Lipschitz es uniformemente continua, y toda aplicación uniformemente continua es continua.
\end{proposicion}
\begin{proof}
Si $L>0$, para una tolerancia $\eps>0$ tomamos $\delta=\eps/L$. Si $L=0$, todas las imágenes coinciden y sirve cualquier $\delta>0$. En ambos casos se cumple la definición. Para la segunda afirmación fijamos $x\in X$: el mismo $\delta$ que funciona para todos los pares funciona, en particular, para los pares $(x,y)$. Ésta es precisamente la continuidad en $x$.
\end{proof}

\begin{ejemplo}[Qué significa aquí discreto]
Si $d(x,y)=1$ para $x\ne y$, toda aplicación de $X$ a cualquier espacio métrico es uniformemente continua: con $\delta=1$ la condición $d(x,y)<\delta$ obliga a $x=y$. El argumento vale siempre que exista $c>0$ con $d(x,y)\geq c$ para puntos distintos; se dice entonces que $X$ es uniformemente discreto. Tener solamente topología discreta no basta. En $X=\{1/n:n\in\N\}\subset\R$, todos los puntos son aislados, pero $f(1/n)=n$ no es uniformemente continua: las distancias entre $1/n$ y $1/(n+1)$ tienden a cero, mientras que las imágenes distan $1$.
\end{ejemplo}

\begin{ejemplo}[Continuidad sin uniformidad]\label{ch07:cuadrado}
La función $f(x)=x^2$ es continua en $\R$, pero no uniformemente continua. Tomemos $x_n=n$ y $y_n=n+1/n$. Entonces
\[
 |x_n-y_n|=1/n\longrightarrow0,
 \qquad |f(x_n)-f(y_n)|=2+1/n^2\geq2.
\]
Si existiera un $\delta$ uniforme para $\eps=1$, un $n$ con $1/n<\delta$ contradiría la segunda desigualdad. Este razonamiento fija primero la tolerancia que fracasa y después produce pares arbitrariamente próximos. No es válido deducir una cota inferior proporcional a $\delta$ a partir de la sola desigualdad $|x-y|<\delta$.
\end{ejemplo}

\begin{proposicion}[Composición]\label{ch07:composicion}
Si $f:X\to Y$ y $g:Y\to Z$ son uniformemente continuas, también lo es $g\circ f$.
\end{proposicion}
\begin{proof}
Fijado $\eps>0$, la uniformidad de $g$ proporciona $\eta>0$ tal que $\rho(u,v)<\eta$ implica $\sigma(g(u),g(v))<\eps$. Aplicamos ahora la uniformidad de $f$ con tolerancia $\eta$: existe $\delta>0$ tal que $d(x,y)<\delta$ implica $\rho(f(x),f(y))<\eta$. Las dos implicaciones encadenadas dan el resultado. El orden de elección es $\eps$, después $\eta$ y finalmente $\delta$.
\end{proof}

\section{Sucesiones cuyos términos terminan estando próximos}
\begin{definicion}\label{ch07:cauchy}
Una sucesión $(x_n)$ de un espacio métrico es \emph{de Cauchy} si
\[
 \forall\eps>0\ \exists N\in\N\ \forall n,m\geq N,
 \qquad d(x_n,x_m)<\eps.
\]
\end{definicion}
\index{Cauchy, sucesion de@Cauchy, sucesión de}
La lectura original utiliza la condición $d(x_n,x_N)<\eps$ para $n\geq N$. Ambas formulaciones son equivalentes: la definición anterior implica esa condición al tomar $m=N$; a la inversa, aplicamos la condición con $\eps/2$ y usamos
\[
 d(x_n,x_m)\leq d(x_n,x_N)+d(x_N,x_m)<\eps.
\]
La condición de Cauchy compara términos de la sucesión entre sí; no presupone que su posible límite pertenezca al espacio.

\begin{teorema}\label{ch07:cauchy-imagen}
Las aplicaciones uniformemente continuas transforman sucesiones de Cauchy en sucesiones de Cauchy.
\end{teorema}
\begin{proof}
Sea $\eps>0$ y elijamos $\delta>0$ por la continuidad uniforme de $f$. Como $(x_n)$ es de Cauchy, existe $N$ tal que $n,m\geq N$ implica $d(x_n,x_m)<\delta$. Por tanto $\rho(f(x_n),f(x_m))<\eps$ para esos índices, que es la condición de Cauchy en $Y$.
\end{proof}
La continuidad ordinaria no basta: $x_n=1/n$ es de Cauchy en $(0,1]$, pero la función continua $x\mapsto1/x$ la transforma en $(n)$, que no es de Cauchy.

\section{Ejercicios resueltos}
\begin{ejercicio}[Lectura 7, ejercicio 1]
Demuestre que la composición de dos funciones Lipschitz es Lipschitz y determine una constante para la composición.
\end{ejercicio}
\begin{solucion}
Si $f:X\to Y$ tiene constante $L$ y $g:Y\to Z$ tiene constante $M$, entonces
\[
 \sigma(g(f(x)),g(f(y)))\leq M\rho(f(x),f(y))\leq MLd(x,y).
\]
Así, $ML$ es una constante válida, incluso cuando alguno de los factores es cero. No necesariamente es la mejor: si $g$ es constante en $f(X)$, la composición tiene constante óptima cero aunque $L,M$ sean positivos.
\end{solucion}

\begin{ejercicio}[Lectura 7, ejercicio 2]
Sean $f:X\to Y$ y $g:Y\to Z$ continuas. Determine si $g\circ f$ debe ser uniformemente continua cuando: \begin{enumerate}[label=(\roman*)]
\item $f$ es uniformemente continua;
\item $g$ es uniformemente continua.
\end{enumerate}
\end{ejercicio}
\begin{solucion}
En ambos casos la respuesta es negativa. En (i) tomamos $X=Y=Z=\R$, $f(x)=x$ y $g(x)=x^2$. La identidad es $1$-Lipschitz y la composición es el cuadrado, que falla por el ejemplo~\ref{ch07:cuadrado}. En (ii) intercambiamos los papeles: $f(x)=x^2$ y $g(x)=x$. Ahora $g$ es $1$-Lipschitz y la composición vuelve a ser $x^2$.
\end{solucion}

\begin{ejercicio}[Lectura 7, ejercicio 3]
Sea $A$ un cerrado no vacío de un espacio métrico $X$. Pruebe que $x\mapsto\dist(x,A)$ es uniformemente continua.
\end{ejercicio}
\begin{solucion}
Para $a\in A$, la desigualdad triangular da $d(x,a)\leq d(x,y)+d(y,a)$. Tomando ínfimos en $a$ obtenemos
\[
 \dist(x,A)\leq d(x,y)+\dist(y,A).
\]
Al intercambiar $x,y$ concluimos
\[
 |\dist(x,A)-\dist(y,A)|\leq d(x,y).
\]
La función es $1$-Lipschitz. La prueba no exige que $A$ sea cerrado; esa hipótesis solamente permite identificar $\dist(x,A)=0$ con $x\in A$. Se agrega la condición $A\ne\varnothing$ porque la distancia al vacío no es una función real finita.
\end{solucion}

\begin{ejercicio}[Lectura 7, ejercicio 4]
Dé una función uniformemente continua entre espacios métricos que no sea abierta.
\end{ejercicio}
\begin{solucion}
La función constante $f:\R\to\R$, $f(x)=0$, es $0$-Lipschitz. Sin embargo, la imagen del abierto $(-1,1)$ es $\{0\}$, que no es abierto en $\R$. La continuidad controla preimágenes de abiertos; ser abierta controla imágenes y es una condición distinta.
\end{solucion}

\begin{ejercicio}[Lectura 7, ejercicio 5]\label{ch07:ej-extension}
Sea $A$ abierto en un espacio métrico $X$ y sea $f:X\to\R$ uniformemente continua en $A$. Demuestre que $f|_A$ se extiende continuamente a $\cl A$. Dé un contraejemplo si se reemplaza $\R$ por un espacio métrico arbitrario.
\end{ejercicio}
\begin{solucion}
Probaremos la afirmación más fuerte en que $A$ es cualquier subconjunto y el espacio de llegada $Y$ es completo. Si $A=\varnothing$, la extensión tiene dominio vacío y no hay nada que definir. Para $x\in\cl A$ elegimos $a_n\in A$ con $d(a_n,x)<1/n$. Esta sucesión converge y es de Cauchy. Por el teorema~\ref{ch07:cauchy-imagen}, $(f(a_n))$ es de Cauchy; la completitud de $Y$ proporciona un límite. Definimos
\[
 g(x)=\lim_{n\to\infty}f(a_n).
\]
La elección no afecta al valor: si $b_n\to x$ en $A$, entonces $d(a_n,b_n)\leq d(a_n,x)+d(x,b_n)\to0$; la uniformidad implica $\rho(f(a_n),f(b_n))\to0$. Tomando límites, los dos valores propuestos coinciden. Si $x\in A$, podemos elegir la sucesión constante $a_n=x$, por lo que $g(x)=f(x)$.

Además, $g$ es uniformemente continua. Dado $\eps>0$, elijamos $\delta>0$ de modo que $d(a,b)<\delta$ implique $\rho(f(a),f(b))<\eps/2$ para $a,b\in A$. Si $x,y\in\cl A$ y $d(x,y)<\delta/3$, elegimos $a_n\to x$, $b_n\to y$. Para $n$ suficientemente grande, $d(a_n,b_n)<\delta$; al tomar límites queda $\rho(g(x),g(y))\leq\eps/2<\eps$. La extensión es única porque dos aplicaciones continuas que coinciden en un conjunto denso coinciden en su clausura: evaluamos ambas sobre una sucesión $a_n\to x$.

Para el contraejemplo tomamos $X=\R$, $A=(0,1)$ y $Y=(0,1)$ con la distancia usual. Definimos $f(x)=x$ en $A$ y $f(x)=1/2$ fuera de $A$. Su restricción a $A$ es $1$-Lipschitz. Una extensión continua $g:[0,1]\to(0,1)$ tendría $g(0)=\lim_{n\to\infty}g(1/(n+2))=0$, imposible en el espacio de llegada. Lo que falla es su completitud.
\end{solucion}

\begin{ejercicio}[Lectura 7, ejercicio 6]
Para subconjuntos no vacíos $A,B$ de $X$, defina $d(A,B)=\inf\{d(a,b):a\in A,b\in B\}$. Pruebe que $f:X\to Y$ es uniformemente continua si y sólo si
\[
 d(A,B)=0\Longrightarrow\rho(f(A),f(B))=0
\]
para todos esos conjuntos.
\end{ejercicio}
\begin{solucion}
Si $f$ es uniformemente continua, dado $\eps>0$ elegimos el correspondiente $\delta>0$. La igualdad $d(A,B)=0$ permite encontrar $a\in A,b\in B$ con $d(a,b)<\delta$. Entonces $\rho(f(A),f(B))\leq\rho(f(a),f(b))<\eps$. Un número no negativo menor que todo $\eps>0$ es cero.

Para el recíproco hay que cuidar las \emph{distancias cruzadas}: saber que $\rho(f(x_n),f(y_n))\geq\eps_0$ no basta por sí solo para separar los conjuntos de imágenes. Si $f$ no es uniformemente continua, existen $\eps_0>0$ y pares $x_n,y_n$ tales que
\[
 d(x_n,y_n)<1/n,\qquad \rho(p_n,q_n)\geq\eps_0,
 \quad p_n=f(x_n),\ q_n=f(y_n).
\]
Pongamos $r=\eps_0/8$. Si alguna bola $B(z,r)$ contiene infinitos $p_n$, retenemos esos índices $I$. Para $i,j\in I$, tenemos $\rho(q_j,z)\geq\eps_0-r$ y, por tanto, $\rho(p_i,q_j)>\eps_0-2r$. Si alguna bola contiene infinitos $q_n$, el argumento con los papeles intercambiados da la misma separación.

Queda el caso en que toda bola de radio $r$ contiene sólo un número finito de términos de cada sucesión, contando sus índices. Elegimos inductivamente índices crecientes $n_k$. Una vez elegidos los primeros $k$, descartamos todos los índices $n$ para los que $p_n\in B(q_{n_j},r)$ o $q_n\in B(p_{n_j},r)$ para algún $j\leq k$. Sólo descartamos una unión finita de conjuntos finitos; por ello queda un índice mayor que $n_k$. La sucesión retenida satisface $\rho(p_{n_i},q_{n_j})\geq r$ si $i\ne j$, y la misma desigualdad vale cuando $i=j$ porque $\eps_0\geq r$.

En cualquiera de los casos obtenemos un conjunto infinito de índices $I$ y una constante $c>0$ tales que $\rho(p_i,q_j)\geq c$ para todo $i,j\in I$. Para $A=\{x_i:i\in I\}$ y $B=\{y_i:i\in I\}$ se cumple $d(A,B)=0$, pues $I$ no está acotado y $d(x_i,y_i)<1/i$, pero $\rho(f(A),f(B))\geq c$. Esto contradice la propiedad supuesta.
\end{solucion}

\begin{ejercicio}[Lectura 7, ejercicio 7]
Suponga que $f:X\to Y$ satisface $\rho(f(x),f(y))\geq\alpha d(x,y)$ para alguna $\alpha>0$. Pruebe que es inyectiva, que su inversa sobre la imagen es Lipschitz y que $f$ es abierta si $f(X)$ es abierto en $Y$.
\end{ejercicio}
\begin{solucion}
Si $f(x)=f(y)$, la desigualdad da $0\geq\alpha d(x,y)$ y por tanto $x=y$. Para $u=f(x),v=f(y)$ obtenemos
\[
 d(f^{-1}(u),f^{-1}(v))\leq\alpha^{-1}\rho(u,v).
\]
Así, $f^{-1}:f(X)\to X$ es continua. Si $U\subset X$ es abierto, $f(U)=(f^{-1})^{-1}(U)$ es abierto en el subespacio $f(X)$. Cuando $f(X)$ es abierto en $Y$, todo abierto suyo es abierto en $Y$, y $f$ es una aplicación abierta. La hipótesis no afirma que $f$ sea continua; controla únicamente su inversa.
\end{solucion}

\section{Qué aporta la uniformidad}
La misma estimación sirve para todos los puntos, preserva la condición de Cauchy y permite extender funciones a puntos límite cuando el espacio de llegada contiene los límites necesarios. El siguiente capítulo estudia precisamente esa propiedad del espacio de llegada: la completitud.
